% Chinese Remainder Theorem (CRT)

Solves a system of simultaneous modular congruences $x \equiv r_i \pmod{m_i}$ for pairwise coprime moduli in $O(N \log(\max M_i))$ time complexity using the Extended Euclidean Algorithm.

\begin{lstlisting}
long long extgcd(long long a, long long b, long long &x, long long &y) {
    if (b == 0) { x = 1; y = 0; return a; }
    long long x1, y1;
    long long d = extgcd(b, a % b, x1, y1);
    x = y1;
    y = x1 - y1 * (a / b);
    return d;
}

// Solves a system of 2 congruences: x = r1 (mod m1) and x = r2 (mod m2).
// Works even if m1 and m2 are NOT coprime. Returns {-1, -1} if no solution exists.
pair<long long, long long> crt_two(long long r1, long long m1, long long r2, long long m2) {
    long long x, y;
    long long g = extgcd(m1, m2, x, y);
    if ((r2 - r1) % g != 0) return {-1, -1};
    
    long long lcm = m1 / g * m2;
    long long mod = m2 / g;
    long long x0 = ((x % mod) + mod) % mod;
    long long factor = (r2 - r1) / g;
    long long ans = (r1 + (__int128)factor * x0 % lcm * m1) % lcm;
    if (ans < 0) ans += lcm;
    return {ans, lcm};
}

// Solves a system of N congruences. Returns {-1, -1} if no solution exists.
// Input format: vectors of pairs {remainder, modulus}
pair<long long, long long> crt(const vector<pair<long long, long long>>& equations) {
    if (equations.empty()) return {0, 1};
    long long r_ans = equations[0].first;
    long long m_ans = equations[0].second;
    
    for (size_t i = 1; i < equations.size(); i++) {
        auto [r, m] = crt_two(r_ans, m_ans, equations[i].first, equations[i].second);
        if (m == -1) return {-1, -1};
        r_ans = r;
        m_ans = m;
    }
    return {r_ans, m_ans};
}
\end{lstlisting}

To find a number $x$ that satisfies $x \equiv 2 \pmod 3$, $x \equiv 3 \pmod 5$, and $x \equiv 2 \pmod 7$:
\begin{lstlisting}
vector<pair<long long, long long>> eq = {{2, 3}, {3, 5}, {2, 7}};
pair<long long, long long> ans = crt(eq);
if (ans.second != -1) {
    long long x = ans.first;   // Unique solution modulo LCM (23)
    long long lcm = ans.second; // LCM of all moduli (105)
}
\end{lstlisting}
